\section{Proof by Contradiction: Power and Danger}
\chapterline

Proof by contradiction is one of the most powerful and psychologically dramatic tools in mathematics. When direct attacks stall and contrapositive restructuring offers no purchase, contradiction allows you to assume the exact opposite of what you hope to prove, step inside that hypothetical counter-world, and walk forward until logical reality shatters.

The English mathematician G.~H.~Hardy famously described it in \emph{A Mathematician's Apology}:
\begin{quote}
``Reductio ad absurdum, which Euclid loved so much, is one of a mathematician's finest weapons. It is a far finer gambit than any chess gambit: a chess player may offer the sacrifice of a pawn or even a piece, but a mathematician offers the game.''
\end{quote}

Yet because of its dramatic power, beginners frequently misuse it. Contradiction is not an excuse to generate vague fog until something ``feels wrong.'' A contradiction proof succeeds only when two crystal-clear, mutually exclusive mathematical facts collide head-on: $A \land \neg A$.

\begin{historybox}
Around 500 BCE, the followers of Pythagoras lived by a mystical doctrine: ``All is Number.'' To the Pythagoreans, every relationship in nature, music, and geometry could be expressed as a ratio of whole numbers. When studying a right triangle with legs of length $1$, they sought two whole numbers $a$ and $b$ whose ratio $a/b$ measured the hypotenuse.

According to ancient legend, a member named Hippasus of Metapontum realized that assuming such a fraction exists leads to an inescapable trap: the numerator $a$ must be even, which forces the denominator $b$ to be even, destroying the assumption that the fraction was reduced to lowest terms. The diagonal was incommensurable with the side. Legend has it that the brotherhood was so unsettled by this scandal that they drowned Hippasus at sea to bury the secret. Yet two centuries later, Euclid recorded this exact parity contradiction in Book X of the \emph{Elements}, alongside his legendary proof that prime numbers never run out.

In the early twentieth century, the Dutch mathematician L.E.J. Brouwer criticized proof by contradiction for being non-constructive. Brouwer argued that proving an object exists by showing that its absence causes a contradiction leaves you empty-handed: it proves non-existence is impossible, but offers no recipe or algorithm to construct the object. In modern computer science, constructive proofs are prized because algorithms require tangible witnesses. Yet proof by contradiction remains our premier weapon for establishing impossibility results: proving that no fraction can equal $\sqrt{2}$, and no finite list contains all primes.
\end{historybox}

\subsection{Visualizing the split reality}

\begin{center}
\begin{tikzpicture}[
    font=\small,
    rootnode/.style={draw=gray!70, fill=gray!8, rounded corners=3pt, thick, text width=4.5cm, align=center, inner sep=6pt},
    realnode/.style={draw=targetgreen!80, fill=targetgreen!10, rounded corners=3pt, thick, text width=4.0cm, align=center, inner sep=6pt},
    counterbox/.style={draw=trapred!70, fill=trapred!4, rounded corners=4pt, dashed, thick, inner sep=8pt},
    cnode/.style={draw=trapred!70, fill=white, rounded corners=2pt, text width=5.0cm, align=center, inner sep=4pt}
]

% Root Node
\node[rootnode] (root) at (2.2, 4.4) {\textbf{Standard Mathematical Universe}\\Axioms, Definitions, Known Theorems};

% Real Branch (Left)
\node[realnode] (real) at (-2.8, 1.3) {
    \textbf{\color{targetgreen!90!black}True Reality}\\[0.3em]
    \scriptsize The desired statement $P$ holds.\\Consistent deduction continues.
};
\draw[-{Latex[length=2.5mm]}, thick, targetgreen!80] (0.5, 3.9) -- (-2.0, 2.2);

% Counter-World Box (Right)
\node[counterbox, minimum width=6.2cm, minimum height=4.6cm] (cbox) at (6.0, 1.3) {};
\node[anchor=north, font=\scriptsize\scshape, text=trapred!90!black] at (6.0, 3.5) {Hypothetical Counter-World: Assume $\neg P$};

\draw[-{Latex[length=2.5mm]}, thick, trapred!80, dashed] (3.8, 3.9) -- (5.5, 3.6);

% Nodes inside Counter-World
\node[cnode] (step1) at (6.0, 2.4) {\footnotesize \textbf{Assume:} $\sqrt{2} = a/b$ in lowest terms\\($\gcd(a, b) = 1$)};
\node[cnode] (step2) at (6.0, 1.3) {\footnotesize \textbf{Deduce:} $a^2 = 2b^2 \implies a$ is even\\$\implies 4k^2 = 2b^2 \implies b$ is even};

% Collision Node
\node[draw=trapred!90, fill=trapred!15, rounded corners=3pt, thick, text width=5.2cm, align=center, inner sep=5pt] (rupture) at (6.0, -0.1) {
    \textbf{\color{trapred!90!black}LOGICAL FRACTURE: $A \land \neg A$}\\[0.2em]
    \scriptsize Both $a$ and $b$ are even ($\gcd(a, b) \ge 2$)\\
    \textbf{Direct collision with} $\gcd(a, b) = 1$!
};

\draw[-{Latex[length=2mm]}, thick, trapred!80] (step1) -- (step2);
\draw[-{Latex[length=2mm]}, thick, trapred!80] (step2) -- (rupture);

% Elimination arrow pointing back to Real
\draw[-{Latex[length=3mm]}, very thick, quotientorange] (cbox.west) -- node[above=3pt, font=\scriptsize\bfseries, text=quotientorange]{Counter-world collapses} node[below=3pt, font=\scriptsize, text=quotientorange]{$P$ is established} (real.east);

\end{tikzpicture}
\end{center}

\subsection{\texorpdfstring{The irrationality of $\sqrt{2}$}{The irrationality of sqrt(2)}}

The classical proof of the irrationality of $\sqrt{2}$ is preserved in Book X of Euclid's \emph{Elements}. It was one of the great intellectual shocks of antiquity: geometry produced lengths (such as the diagonal of a unit square) that arithmetic, if restricted to ratios of integers, could not measure.

\begin{theorem}
$\sqrt{2}$ is irrational.
\end{theorem}

\begin{proof}
Suppose, for contradiction, that $\sqrt{2}$ is rational. 

By definition of rational numbers, we may express $\sqrt{2}$ as a ratio of integers:
\[
\sqrt{2} = \frac{a}{b}, \qquad \text{where } a, b \in \Z \text{ and } b \neq 0.
\]
Furthermore, we can choose this fraction in \textbf{lowest terms}, meaning that $a$ and $b$ share no common factors other than $1$ (that is, their greatest common divisor $\gcd(a, b) = 1$). In particular, $a$ and $b$ cannot both be even.

Now we perform lawful algebraic operations:
\begin{align*}
\sqrt{2} &= \frac{a}{b} \\
2 &= \frac{a^2}{b^2} \qquad \text{(squaring both sides)} \\
a^2 &= 2b^2 \qquad \text{(multiplying by } b^2\text{)}.
\end{align*}
Since $b$ is an integer, $b^2$ is an integer. Thus $a^2$ is an integer multiple of $2$, which means $a^2$ is even.

We proved in Chapter 7 that for any integer $a$, if $a^2$ is even, then $a$ must be even. Therefore:
\[
a = 2k \quad \text{for some integer } k.
\]

Substitute this expression for $a$ back into our equation $a^2 = 2b^2$:
\begin{align*}
(2k)^2 &= 2b^2 \\
4k^2 &= 2b^2 \\
2k^2 &= b^2 \qquad \text{(dividing both sides by } 2\text{)}.
\end{align*}
Because $k$ is an integer, $k^2$ is an integer. This equation shows that $b^2$ is an integer multiple of $2$, meaning $b^2$ is even. Again invoking our theorem, since $b^2$ is even, $b$ must be even!

Now look at what has happened:
\begin{itemize}
\item We proved that $a$ is even ($2 \mid a$).
\item We proved that $b$ is even ($2 \mid b$).
\end{itemize}
This means $2$ is a common divisor of $a$ and $b$. But this directly contradicts our initial specification that $a/b$ was chosen in lowest terms ($\gcd(a, b) = 1$).

Both statements cannot be true simultaneously in mathematics. Because every deductive step followed lawful arithmetic rules, the initial assumption that $\sqrt{2}$ can be written as a ratio of integers must be false. 

Therefore, $\sqrt{2}$ is irrational.
\end{proof}

Notice how clean the trap is: it is a parity ping-pong match. Assuming rationality forces $a$ to be even, which forces $b$ to be even, which destroys the fraction's lowest-terms certificate.

\subsection{The infinitude of primes}

Another milestone of contradiction is Euclid's proof that prime numbers never run out.

\begin{theorem}[Euclid]
There are infinitely many prime numbers.
\end{theorem}

\begin{proof}
Suppose, for contradiction, that the set of all prime numbers is finite. 

Because the set is finite, we can write down all prime numbers in a complete, terminating list:
\[
\{p_1, p_2, p_3, \ldots, p_n\}.
\]
Here $p_1 = 2$, $p_2 = 3$, $p_3 = 5$, up to the supposed final prime $p_n$.

Now construct the integer $N$ obtained by multiplying all of these primes together and adding $1$:
\[
N = (p_1 \cdot p_2 \cdot p_3 \cdots p_n) + 1.
\]
Because $p_1 = 2 \ge 2$, the number $N$ is strictly greater than $1$. Every integer greater than $1$ has at least one prime divisor. Therefore, $N$ must be divisible by some prime number $q$.

Is this prime divisor $q$ on our list of all primes?
\begin{itemize}
\item If $q = p_i$ for some index $i \in \{1, 2, \ldots, n\}$, then $q$ divides the product $P = p_1 p_2 \cdots p_n$.
\item But $q$ also divides $N = P + 1$.
\item If an integer divides two numbers, it must divide their difference:
\[
q \mid (N - P) \implies q \mid 1.
\]
\end{itemize}
The only positive divisor of $1$ is $1$. But by definition, a prime number must be greater than $1$, so $q \ge 2$. A number greater than or equal to $2$ cannot divide $1$.

This is an impossibility. The assumption that the list $\{p_1, \ldots, p_n\}$ contained all prime numbers has collapsed. Therefore, the set of prime numbers is infinite.
\end{proof}

\subsection{The common trap: phantom contradictions}

Beginners sometimes attempt a proof by contradiction when a direct proof is sitting right in front of them:
\begin{quote}
\itshape ``Theorem: If $n$ is even, then $n+1$ is odd.\\
Attempted contradiction: Suppose $n$ is even and $n+1$ is even. Then $n = 2k$ and $n+1 = 2m$. Subtracting gives $1 = 2(m-k)$, which means $1$ is even. Contradiction!''
\end{quote}
While logically valid, this is stylistic clutter. You do not need to assume $n+1$ is even to show that $n+1 = 2k+1$. Contradiction should be reserved for cases where denying the conclusion gives you an essential structural foothold that direct calculation cannot reach.

\begin{trapbox}
Always name the exact collision in your final sentence. Never end a proof with ``\ldots which is impossible, so the theorem is proved.'' State plainly: ``This contradicts the fact that $q \ge 2$'' or ``This contradicts our assumption that $\gcd(a, b) = 1$.''
\end{trapbox}

\subsection{What contradiction unlocked---and why we need counterexamples}

Proof by contradiction unlocked an extraordinary capability that direct calculations can never match: \textbf{proving impossibility}. A direct calculation can only confirm properties of objects that exist. Contradiction traps a hypothetical non-existent object in a logical vice, showing that its existence would fracture arithmetic itself. Without contradiction, we could never prove that $\sqrt{2}$ cannot be written as a fraction, that prime numbers never terminate, or that there is no general formula for fifth-degree polynomials.

Yet proof by contradiction has an honest limitation: \textbf{it cannot tell you why an exploratory attempt fails}. If you suspect a conjecture is true and try to prove it by assuming its negation, you may spend weeks searching for a contradiction that never appears. Did your proof stall because you missed a clever algebraic identity, or because the conjecture is fundamentally false? 

To answer that question, mathematicians do not wait passively for a contradiction; they go on the offensive. In Chapter~9 (\emph{Examples, Non-Examples, and Counterexamples}), we learn the art of empirical stress-testing: finding minimal objects that break false claims and discovering how mathematical theories develop through proofs and refutations.

\subsection{Exercises}

\begin{enumerate}
\item Prove that $\sqrt{3}$ is irrational. (You may use without proof the lemma: for any integer $n$, if $3 \mid n^2$, then $3 \mid n$.)
\item Prove that $\log_{10}(2)$ is irrational. (Hint: Assume $\log_{10}(2) = a/b$ with positive integers $a, b$, rewrite in exponential form $10^a = 2^b$, and consider prime factorization.)
\item Prove by contradiction that there is no smallest positive rational number.
\item Prove that the sum of a rational number and an irrational number is irrational.
\item Re-read Euclid's proof of the infinitude of primes. Does the proof claim that $N = p_1 p_2 \cdots p_n + 1$ is always a prime number? Check the value of $(2 \cdot 3 \cdot 5 \cdot 7 \cdot 11 \cdot 13) + 1 = 30031$. Is $30031$ prime? (Compute $59 \times 509$.) Explain why Euclid's proof remains completely sound regardless.
\end{enumerate}

\newpage
